\section{Related Work}
\label{sec:relatedwork}
\textbf{Mechanisms for Secure Speculation}
Many academic works \cite{ainsworth2020muontrap, cheang2019formal, canella2019systematic, barber2019specshield, fustos2019spectreguard, yan2018invisispec, yu2019speculative, saileshwar2019cleanupspec, weisse2019nda, khasawneh2019safespec, loughlin2021dolma, ainsworth2021ghostminion} present hardware mechanisms or hardware-software co-design to defend against these attacks.
They sometimes aim for a stronger security property than us (speculative constant time) but do so within threat models that restrict adversarial power to control or observe only a subset of microarchitectural state (e.g., only consider cache side channels) and only consider specific forms of speculation (e.g., branch predictors cannot be updated speculatively).
These threat models could be seen as indirectly building assumptions on the microarchitectural isolation and the speculation model, making these mechanisms potential add-ons to a Citadel that satisfies these isolation requirements.
As a result, we see our work as complementary. We note that very few of these mechanisms have been implemented in hardware, i.e., FPGAs or silicon.
Many attacks that break these defense mechanisms rely on exploiting details that are often glossed over in simulation (like not faithfully simulating the memory hierarchy, for instance).

\noindent\textbf{Enclave Platforms}
Sanctum \cite{costan2016sanctum} introduces the concept of a privileged security monitor by decoupling enclave management functionalities from hardware mechanisms. 
Our platform draws many concepts from Sanctum such as the dual page table and the memory isolation mechanism. Sanctum also considers some cache side channels in its threat model.
Nonetheless, Sanctum is not concerned with speculative side channels and no functioning enclave system was ever built.
Other platforms such as Keystone \cite{lee2020keystone}, CURE \cite{bahmani2021cure}, or \cite{enclavecaches} address cache side channels, but ignore other side channels and only consider in-order processors. 
Many other academic platforms exist \cite{ferraiuolo2017komodo, feng2021scalable, weiser2019timber, mccune2008flicker, azab2011sice, xia2013architecture, deng2014equalvisor, brasser2019sanctuary, sun2015trustice, koeberl2014trustlite, aegis, xom}, but none of them consider speculative side channels or transient execution attacks to the extent of Citadel. 

\noindent\textbf{MI6}
Citadel builds on the open-source code provided in MI6 \cite{bourgeat2019mi6} and reuses many of its hardware mechanisms.
The MI6 paper focused on evaluating the \emph{performance} of these mechanisms on an FPGA
by running SPEC benchmarks directly on top of (insecure) Linux and using Linux system calls as hooks to project the cost of a potential context switch to the SM.
The platform (likely missing the necessary software) did not evaluate a full enclave scenario (neither did it guarantee any security when running SPEC).
In this paper, on top of implementing secure shared-memory and several mechanisms to improve performance and configurability, we have built and evaluated an end-to-end hardware-software platform, focusing on enclaves that run security-aware workloads and uncovering practical insights for enclave platforms.
Note that Citadel's performance numbers for SPEC would be similar to MI6's evaluation as Citadel's LLC can be configured to match the partition size used in \cite{bourgeat2019mi6}.

\noindent\textbf{Enclave Execution Models}
Recent works have tried to make memory sharing and enclave execution more extensible and configurable \cite{van2023cheri, li2021confidential, lee2022cerberus, park2020nested, steinegger2021servas}.
Those works are focusing on architectural capabilities and isolation, but ignore security challenges introduced by microarchitectural side-channels and transient execution attacks.
Under such a threat model, even basic communication primitives like Citadel public shared memory with the OS are not straightforward to implement efficiently and securely.